\section{Axis 2: Planner-Content Trust Is Capacity-Dependent in Opposite Directions Across Planners}
\label{sec:axis2}

The cleanest signal in the v2-vs-v3 comparison
(Table~\ref{tab:axis1_v2_v3}) is the inversion of the
\Qwen{0.5B}-vs-\Qwen{1.5B} planner split.  At v2 capacity ($4/7$
modules, $10$M parameters), the larger planner helps and the
smaller one hurts: \cthree{}\code{p Q-1.5B} reaches $67.0\%$ versus
\cthree{}\code{p Q-0.5B} at $60.0\%$ ($\Delta=+7$\,pp).  At v3
capacity ($7/7$ modules, $22$M parameters), the directions invert:
\cthree{}\code{p Q-1.5B} drops to $54.0\%$ while
\cthree{}\code{p Q-0.5B} rises to $65.0\%$ ($\Delta=-11$\,pp).
This is a previously-unmeasured axis: LoRA capacity differentially
affects how much the model \emph{trusts} the content of an
upstream plan, and the direction of the effect is opposite for the
two planners.  We have two capacity rungs (v2 at $10$M and v3 at
$22$M trainable parameters); the data supports a directional split
across planners but does not yet trace a fitted non-monotonic
curve, and we are careful not to overclaim that here.

\subsection{The Pattern}
\label{subsec:axis2_pattern}

Table~\ref{tab:axis2_planner_split} extracts the planner conditions
from \S\ref{sec:axis1} and adds the v2 $\to$ v3 deltas.

\begin{table}[h]
\centering
\caption{Planner-content trust as a function of LoRA capacity.
At v2 the bigger planner helps; at v3 it catastrophically regresses
while the smaller planner improves.  Static-prefix conditions
(\code{C2hint}, \code{C2empty}) are unaffected by the capacity bump
and are listed for contrast.}
\label{tab:axis2_planner_split}
\smallskip
\begin{tabular}{lrrrr}
\toprule
Condition & Base & v2 & v3 & v2 $\to$ v3 \\
\midrule
\code{C2hint} (static) & $68\%$ & $73.5\%$ & $73.5\%$ & $0.0$ \\
\code{C2empty} (static) & $66\%$ & $73.0\%$ & $74.0\%$ & $+1.0$ \\
\cthree{}\code{p Q-0.5B} (content) & $64\%$ & $60.0\%$ & $\mathbf{65.0\%}$ & $\mathbf{+5.0}$ \\
\cthree{}\code{p Q-1.5B} (content) & $60\%$ & $67.0\%$ & $\mathbf{54.0\%}$ & $\mathbf{-13.0}$ \\
\bottomrule
\end{tabular}
\end{table}

The capacity bump is invisible on the static-prefix conditions
(both within $1$\,pp of v2) and amplifies in opposite directions
on the content-rich conditions.  Since the only thing that
distinguishes \cthree{}\code{p Q-0.5B} from \cthree{}\code{p Q-1.5B}
is the content of the prefix --- the surface shape is identical ---
the capacity-amplified differential is an effect on plan-content
processing, not on plan shape.

\subsection{Reading}
\label{subsec:axis2_reading}

The natural interpretation: full-FFN coverage at v3 makes the
adapter \emph{more} sensitive to the content of the prefix.  For
\Qwen{0.5B}, the prefixes are simple and often noisy, and the
extra sensitivity is corrective --- the LoRA learns to ignore weak
content effectively, recovering $+5$\,pp.  For \Qwen{1.5B}, the
prefixes are more sophisticated and contain plan-internal patterns
that the adapter latches onto.  The increased sensitivity makes
the model overfit to those patterns at training time and they do
not generalize at evaluation, regressing past the base $60\%$
baseline by $-6$\,pp net.

This reading is consistent with the static-prefix conditions being
unaffected by the capacity bump: \code{"Plan: "} and
\code{"Let's think step by step."} have no content for the adapter
to over-fit to, so additional capacity has no per-content
attractor to drift toward.  Only the planner conditions, where
the prefix carries question-specific content, expose the trust
modulation.

For deployment, the practical implication is that the v2/v3 choice
is not a uniform Pareto frontier: it depends on the size of the
upstream planner.  A pipeline that pairs the prefix-robustness
LoRA with a $\sim 0.5$B planner should ship v3; a pipeline that
pairs it with a $\sim 1.5$B planner should ship v2.  We did not
test \Qwen{7B} planners; the trust modulation may continue or it
may bend back at larger planner scale, and characterizing the full
curve is left to future work.

\paragraph{Status.}
This paper characterizes the axis but does not fix it.  A
content-trust adapter --- LoRA training that varies plan content
quality at fixed surface shape, rather than varying surface shape
at fixed content as our prefix-robust dataset does --- is a
natural follow-up.  We sketch the design in
\S\ref{sec:discussion}.
